\newcommand{\figuraGasIdealMicroscopico}{%
\begin{figure}[htbp]
    \centering
    \begin{tikzpicture}[>=Latex,font=\small]
        \node[font=\bfseries] at (6.4,5.45)
            {Interpretación microscópica del gas ideal};

        \draw[very thick,rounded corners=3pt,fill=blue!4]
            (0.5,0.8) rectangle (5.6,4.7);
        \node[font=\bfseries] at (3.05,4.98) {$T_1$};
        \foreach \x/\y/\dx/\dy in {
            1.10/1.35/0.35/0.20,2.15/1.20/-0.30/0.22,3.15/1.45/0.28/-0.25,
            4.40/1.20/-0.35/-0.15,1.30/2.35/0.25/-0.28,2.55/2.25/-0.32/0.18,
            3.70/2.55/0.30/0.22,4.75/2.35/-0.25/0.30,1.05/3.55/0.32/0.16,
            2.20/3.45/-0.25/-0.30,3.35/3.70/0.34/-0.12,4.55/3.55/-0.30/0.20}
        {
            \shade[ball color=blue!65] (\x,\y) circle (0.09);
            \draw[gray!75,-{Latex[length=2mm]}]
                (\x,\y)--(\x+\dx,\y+\dy);
        }

        \draw[very thick,rounded corners=3pt,fill=red!3]
            (7.2,0.8) rectangle (12.3,4.7);
        \node[font=\bfseries] at (9.75,4.98) {$T_2>T_1$};
        \foreach \x/\y/\dx/\dy in {
            7.80/1.35/0.62/0.36,8.85/1.20/-0.55/0.40,9.85/1.45/0.50/-0.48,
            11.10/1.20/-0.62/-0.28,8.00/2.35/0.48/-0.52,9.25/2.25/-0.60/0.34,
            10.40/2.55/0.58/0.40,11.45/2.35/-0.46/0.56,7.75/3.55/0.58/0.30,
            8.90/3.45/-0.48/-0.56,10.05/3.70/0.64/-0.22,11.25/3.55/-0.56/0.38}
        {
            \shade[ball color=red!65] (\x,\y) circle (0.09);
            \draw[red!70!black,-{Latex[length=2mm]}]
                (\x,\y)--(\x+\dx,\y+\dy);
        }

        \node[align=center] at (3.05,0.18)
            {menor energía cinética\\menor presión};
        \node[align=center] at (9.75,0.18)
            {mayor energía cinética\\mayor presión};
        \node[align=center,fill=white,rounded corners=2pt,
              draw=blue!25,inner sep=5pt] at (6.40,-0.80)
            {$\displaystyle
             \frac{1}{2}m\langle v^2\rangle=\frac{3}{2}k_BT$
             \qquad
             $\displaystyle p=\frac{1}{3}\rho\langle v^2\rangle=\rho RT$};
    \end{tikzpicture}
    \caption{Al aumentar la temperatura, aumenta la energía cinética
    molecular y el intercambio de momento con las paredes. Para densidad
    fija, la presión de un gas ideal es proporcional a la temperatura.}
    \label{fig:gas-ideal-microscopico}
\end{figure}%
}